\begin{notechapter}
\chapter{Case Study: Deep Kalman Filters}\label{ch:deep-kalman-filters}
\chaptersubtitle{Recurrent networks inside a latent-variable model: a backward RNN for inference, a forward RNN for the dynamics}

\begin{chapterintro}
A state-space model explains an observed sequence $x_{1:T}$ by a hidden sequence $z_{1:T}$ that evolves
by itself and produces one observation per step. When the evolution and the observation are
linear and Gaussian, the Kalman filter and smoother compute the posterior of the states
exactly. The deep Kalman filter (DKF) replaces the linear maps by neural networks and loses
this exactness: the posterior must be approximated by an \emph{inference network}, trained
together with the model by maximizing the ELBO. This chapter builds that network from the
conditional independences of the model, read off its graph by d-separation. The answer is
precise: the posterior factor of $z_t$ needs $z_{t-1}$ and a summary of the \emph{future} observations, which
is exactly what a backward RNN provides (\cref{ch:sequence-models}). Making the dynamics depend
on the \emph{past} observations through a forward RNN changes the answer, and the inference
network must change with it. The chapter follows the report and the code of the repository
\texttt{CL-BZH/Deep-Kalman-Filters}. You will learn
\begin{itemize}[nosep]
  \item the linear-Gaussian state-space model and its deep version;
  \item Bayesian networks and d-separation, enough to read conditional independences from a
        graph;
  \item the structure of the DKF posterior and the inference network that matches it:
        backward LSTM, combiner, encoder;
  \item the ELBO of the DKF and its one-trajectory estimator, with closed-form Gaussian KL
        terms;
  \item the autoregressive DKF, whose transition is a density depending on the past through a
        forward RNN, and why its posterior then needs the past as well.
\end{itemize}
\prerequisites{\cref{ch:sequence-models,ch:gated-rnn}; the ELBO, its estimation along one
ancestral trajectory and the reparametrization trick (\cref{em-algorithm:sec:mc-elbo});
ancestral sampling (\cref{monte-carlo:def:ancestral}).}
\end{chapterintro}

\section{State-space models}\label{dkf:sec:ssm}

A \emph{linear-Gaussian state-space model} with inputs $u_t$ is
\[
  z_t = Az_{t-1} + Bu_t + w_t,\quad w_t\sim\mathcal N(0, Q),\qquad
  x_t = Cz_t + v_t,\quad v_t\sim\mathcal N(0, R),
\]
with independent noises. Every variable is a linear function of Gaussian variables, so every
conditional distribution of states given observations is Gaussian, with mean and
covariance given by recursions: the Kalman filter computes $p(z_t\mid x_{1:t})$ (Kalman, \emph{A New Approach
to Linear Filtering and Prediction Problems}, 1960), and the Rauch--Tung--Striebel smoother
computes $p(z_t\mid x_{1:T})$ with a backward pass (Rauch, Tung and Striebel, 1965).

The \emph{deep Kalman filter} (Krishnan, Shalit and Sontag, \emph{Deep Kalman Filters}, 2015) keeps
the structure and replaces the linear maps by neural networks:
\[
  z_t\sim\mathcal N\bigl(\mu^z_\theta(z_{t-1}, u_t),\ \diag\sigma^{z\,2}_\theta(z_{t-1}, u_t)\bigr),\qquad
  x_t\sim\mathcal N\bigl(\mu^x_\theta(z_t),\ \diag\sigma^{x\,2}_\theta(z_t)\bigr),
\]
for $t = 1,\dots,T$, from an initial state $z_0$ (fixed at $0$ or learned); the observation distribution can be
any distribution whose parameters a network outputs. Nonlinear networks destroy the
Gaussianity of the posterior: $p_\theta(z_{1:T}\mid x_{1:T})$ has no closed form, and it is approximated by a
distribution $q_\phi(z_{1:T}\mid x_{1:T})$ computed by an inference network, as in a variational autoencoder.
The question of this chapter is which form $q_\phi$ should have.

\section{Reading independences from a graph}\label{dkf:sec:dseparation}

\begin{definition}[Bayesian network]\label{dkf:def:bn}
A \emph{Bayesian network} is a directed acyclic graph whose nodes are random variables
$v_1,\dots,v_n$, together with a joint distribution of the form
\[
  p(v_1,\dots,v_n) = \prod_{k=1}^np\bigl(v_k\mid\mathrm{pa}(v_k)\bigr),
\]
where $\mathrm{pa}(v_k)$ is the set of parents of $v_k$ in the graph.
\end{definition}

The graph says which variables each factor may depend on. It also implies conditional
independences, which d-separation reads off the graph without any computation. A
\emph{path} between two nodes is a sequence of adjacent nodes, \emph{regardless of the direction of the
edges}: in $A\leftarrow B\to C$ there is a path from $A$ to $C$. Each intermediate node $B$ of a path is of one of three
kinds, according to the two edges of the path that meet at it (\cref{dkf:fig:connections}).

\begin{figure}[ht]
\centering
\begin{tikzpicture}[
  v/.style={draw, circle, minimum size=7mm, inner sep=0pt, font=\small},
  ob/.style={v, fill=gray!30},
  ar/.style={-{Stealth[length=1.8mm]}, thick}]
  \node[v] (a1) at (0,0) {$A$}; \node[v] (b1) at (1.3,0) {$B$}; \node[v] (c1) at (2.6,0) {$C$};
  \draw[ar] (a1) -- (b1); \draw[ar] (b1) -- (c1);
  \node[font=\small, align=center] at (1.3,-0.9) {chain\\[-1pt]{\scriptsize blocked if $B$ is observed}};
  \node[v] (a2) at (4.2,0) {$A$}; \node[v] (b2) at (5.5,0.8) {$B$}; \node[v] (c2) at (6.8,0) {$C$};
  \draw[ar] (b2) -- (a2); \draw[ar] (b2) -- (c2);
  \node[font=\small, align=center] at (5.5,-0.9) {fork\\[-1pt]{\scriptsize blocked if $B$ is observed}};
  \node[v] (a3) at (8.4,0.8) {$A$}; \node[v] (b3) at (9.7,0) {$B$}; \node[v] (c3) at (11.0,0.8) {$C$};
  \draw[ar] (a3) -- (b3); \draw[ar] (c3) -- (b3);
  \node[font=\small, align=center] at (9.7,-0.9) {collider\\[-1pt]{\scriptsize blocked unless $B$ or a}\\[-3pt]{\scriptsize descendant is observed}};
\end{tikzpicture}
\caption{The three kinds of intermediate node on a path.}
\label{dkf:fig:connections}
\end{figure}

\begin{definition}[d-separation]\label{dkf:def:dsep}
Let $\mathcal S$ be a set of observed nodes. A path is \emph{blocked} by $\mathcal S$ if it has an intermediate node $B$
such that either $B$ is a chain or fork node and $B\in\mathcal S$, or $B$ is a collider and neither $B$ nor any of its
descendants is in $\mathcal S$. Otherwise the path is \emph{active}. Two sets of nodes $\mathcal X$ and $\mathcal Y$ are \emph{d-separated}
by $\mathcal S$ if every path between a node of $\mathcal X$ and a node of $\mathcal Y$ is blocked.
\end{definition}

\begin{theorem}[Soundness of d-separation]\label{dkf:thm:dsep}
If $\mathcal X$ and $\mathcal Y$ are d-separated by $\mathcal S$, then $\mathcal X\indep\mathcal Y\mid\mathcal S$ in every distribution that factorizes as in
\cref{dkf:def:bn}. Conversely, if they are not d-separated, there is a distribution that
factorizes over the graph in which they are dependent given $\mathcal S$.
\end{theorem}

The theorem is due to Pearl and his collaborators (Pearl, \emph{Probabilistic Reasoning in
Intelligent Systems}, 1988); a complete treatment is in Koller and Friedman, \emph{Probabilistic
Graphical Models}, 2009, Chapter~3. We use it without proof. The second statement is why a
missing d-separation is read as ``dependent in general'': the graph alone does not
guarantee independence, and for generic parameters it does not hold.

\paragraph{The two mechanisms.} \emph{Chains and forks: screening off.} Let $A$ = ``it rains'',
$B$ = ``the grass is wet'', $C$ = ``someone slips'', $A\to B\to C$. Without information about the grass,
rain makes slipping more likely; once the wetness is known, rain tells nothing more, since all
its influence goes through the grass. In the fork ``fever $\leftarrow$ flu $\to$ cough'', fever and cough are
correlated in the population but independent among patients known to have the flu.
\emph{Colliders: explaining away.} Let $A$ = ``it rained'', $C$ = ``the sprinkler was on'', and $B$ = ``the
grass is wet'', $A\to B\leftarrow C$. Rain and sprinkler are independent causes. Once the grass is known to be
wet, learning that it rained makes the sprinkler \emph{less} likely, because the rain already
explains the wet grass: observing a common effect makes its causes dependent. This
phenomenon, and only this one, is called \emph{explaining away}.

\begin{intuition}
Information flows through a chain or a fork unless the middle node is known, and through
a collider only if the collider, or something it causes, is known. Two variables are
conditionally independent when every route between them is cut by one of these rules.
\end{intuition}

\section{The DKF as a Bayesian network}\label{dkf:sec:model}

The DKF of \cref{dkf:sec:ssm} is the Bayesian network of \cref{dkf:fig:model}: each state $z_t$ has the
parents $z_{t-1}$ and $u_t$, and each observation $x_t$ the single parent $z_t$. The inputs are known, not
random: conditioning on them means that the model is a family of distributions indexed by
$u_{1:T}$, and in the graph they are nodes that are always observed. By \cref{dkf:def:bn}, the
generative model is
\begin{equation}\label{dkf:eq:generative}
  p_\theta(x_{1:T}, z_{1:T}\mid u_{1:T}) = \prod_{t=1}^Tp_\theta(z_t\mid z_{t-1}, u_t)\,p_\theta(x_t\mid z_t),
\end{equation}
where for $t = 1$ the factor $p_\theta(z_1\mid z_0, u_1)$ is the distribution of the first state given the initial value $z_0$.

\begin{figure}[ht]
\centering
\begin{tikzpicture}[
  v/.style={draw, circle, minimum size=8.5mm, inner sep=0pt, font=\small},
  ob/.style={v, fill=gray!30},
  ar/.style={-{Stealth[length=1.8mm]}, thick}]
  \node[v, dashed] (z0) at (0,0) {$z_0$};
  \foreach \t in {1,2,3} {
    \node[ob] (u\t) at (2.2*\t-0.9,1.4) {$u_{\t}$};
    \node[v] (z\t) at (2.2*\t,0) {$z_{\t}$};
    \node[ob] (x\t) at (2.2*\t,-1.5) {$x_{\t}$};
    \draw[ar] (u\t) -- (z\t); \draw[ar] (z\t) -- (x\t);
  }
  \node (zd) at (8.4,0) {$\cdots$};
  \node[ob] (uT) at (9.7,1.4) {$u_T$};
  \node[v] (zT) at (10.6,0) {$z_T$};
  \node[ob] (xT) at (10.6,-1.5) {$x_T$};
  \draw[ar] (uT) -- (zT); \draw[ar] (zT) -- (xT);
  \draw[ar] (z0) -- (z1); \draw[ar] (z1) -- (z2); \draw[ar] (z2) -- (z3); \draw[ar] (z3) -- (zd); \draw[ar] (zd) -- (zT);
\end{tikzpicture}
\caption{The DKF as a Bayesian network. Shaded nodes are observed: the inputs always, the
observations when the posterior is computed. The initial value $z_0$ is a constant.}
\label{dkf:fig:model}
\end{figure}

\section{The structure of the posterior}\label{dkf:sec:posterior}

By the product rule, \emph{any} posterior over a sequence can be written
\[
  p_\theta(z_{1:T}\mid x_{1:T}, u_{1:T}) = \prod_{t=1}^Tp_\theta(z_t\mid z_{1:t-1}, x_{1:T}, u_{1:T}) .
\]
The graph simplifies each factor.

\begin{proposition}[Posterior of the DKF]\label{dkf:prop:posterior}
\[
  p_\theta(z_{1:T}\mid x_{1:T}, u_{1:T}) = \prod_{t=1}^Tp_\theta(z_t\mid z_{t-1},\ x_{t:T},\ u_{t:T}) .
\]
The past states $z_{1:t-2}$, the past observations $x_{1:t-1}$ and the past inputs $u_{1:t-1}$ drop out; the future
observations and the future inputs do not.
\end{proposition}

\begin{proof}
We show that $z_t$ is d-separated from $\mathcal P = \{z_{1:t-2}, x_{1:t-1}, u_{1:t-1}\}$ by $\mathcal S = \{z_{t-1}, x_{t:T}, u_{t:T}\}$; by
\cref{dkf:thm:dsep}, the factor then does not depend on $\mathcal P$. Every $x_s$ and every $u_s$ has a single
neighbour ($z_s$), so it cannot be an intermediate node of a path; every path from $\mathcal P$ to $z_t$
therefore runs through the latent chain. In the graph, the only edge between $\{z_0,\dots,z_{t-1}\}$ and
$\{z_t,\dots,z_T\}$ is $z_{t-1}\to z_t$, so the path reaches $z_t$ through $z_{t-1}$, as an intermediate node entered from $z_{t-2}$,
$x_{t-1}$ or $u_{t-1}$:
\[
  z_{t-2}\to z_{t-1}\to z_t\ \ \text{(chain)},\qquad x_{t-1}\leftarrow z_{t-1}\to z_t\ \ \text{(fork)},\qquad u_{t-1}\to z_{t-1}\to z_t\ \ \text{(chain)}.
\]
(If $\mathcal P$ contains $x_{t-1}$ or $u_{t-1}$ itself, that node is the start of the path, and $z_{t-1}$ is still an
intermediate chain or fork node.) Since $z_{t-1}\in\mathcal S$, every such path is blocked.

The future does not drop out: $z_t\to x_t$ is an edge, which nothing can block, and
$z_t\to z_{t+1}\to x_{t+1}$ is a chain whose middle node $z_{t+1}$ is not observed, hence active. For the inputs,
$z_t\to z_{t+1}\leftarrow u_{t+1}$ has the collider $z_{t+1}$, whose descendant $x_{t+1}$ is observed: the path is active.
\end{proof}

\begin{intuition}
Once the state at time $t-1$ is known, the past observations have nothing more to say about
the state at time $t$: their information has been absorbed by $z_{t-1}$. The future observations still
matter, because they were produced by states that $z_t$ influences. A future input matters by
explaining away: if $x_{t+1}$ is surprising, knowing that $u_{t+1}$ pushed the system explains part of the
surprise and changes what $z_t$ must have been.
\end{intuition}

\paragraph{The inference model.} It is natural to give $q_\phi$ the same structure,
\begin{equation}\label{dkf:eq:q}
  q_\phi(z_{1:T}\mid x_{1:T}) = \prod_{t=1}^Tq_\phi(z_t\mid z_{t-1}, x_{t:T}),
\end{equation}
where from now on there are no inputs, as in the data sets of the repository. Then the
family of $q_\phi$ contains the true posterior when the networks are flexible enough, and the
ELBO can be tight. The first DKF paper experimented with mean-field inference networks
that ignore $z_{t-1}$, conditioning $z_t$ on $x_t$, on $x_{t-1:t+1}$, on $x_{1:t}$ or on $x_{1:T}$. They cannot represent the strong
dependence between successive states of the posterior. The follow-up (Krishnan, Shalit and
Sontag, \emph{Structured Inference Networks for Nonlinear State Space Models}, 2017) builds
networks of the form \cref{dkf:eq:q}, called the deep Kalman smoother, from $z_{t-1}$ and an RNN summary of
$x_{t:T}$; the repository follows this structure.

\section{The inference network}\label{dkf:sec:inference}

The conditional $q_\phi(z_t\mid z_{t-1}, x_{t:T})$ is a Gaussian $\mathcal N(\mu_t, \diag\sigma_t^2)$ whose parameters are computed in three
stages (\cref{dkf:fig:network}).

\paragraph{1. Summarize the future.} A backward LSTM (\cref{sequence-models:sec:flip}) computes, in one
pass over the sequence, $b_1,\dots,b_T$ with $b_t$ a function of $x_{t:T}$. One pass gives every summary the factors
need.

\paragraph{2. Combine with the previous state.} A \emph{combiner} merges $z_{t-1}$ and $b_t$ into one vector.
Krishnan et al.\ (2017) use
\[
  h^{\mathrm{comb}}_t = \tfrac12\bigl(\tanh(Wz_{t-1} + c) + b_t\bigr) .
\]
The repository's linear combiner generalizes the weights:
$h^{\mathrm{comb}}_t = L\bigl(\alpha\tanh(Wz_{t-1} + c) + (1 - \alpha)b_t\bigr)$ with a linear output map $L$ and a weight $\alpha$ that can be learned.
At $t = 1$ the ``previous state'' is the initial value $z_0$, which carries no information about the
sequence; a separate weight $\alpha_{\mathrm{init}}$ (by default $0.1$) is used for that step, so that the first
state is inferred mostly from the observations.

\paragraph{3. Encode and sample.} An MLP maps $h^{\mathrm{comb}}_t$ to $(\mu_t, \log\sigma_t^2)$, and a sample is drawn with the
reparametrization trick (\cref{em-algorithm:sec:mc-elbo}):
$\dot z_t = \mu_t + \sigma_t\odot\varepsilon_t$, $\varepsilon_t\sim\mathcal N(0, I)$. The sample $\dot z_t$ is the ``previous state'' of step $t+1$.

\begin{figure}[ht]
\centering
\begin{tikzpicture}[
  cell/.style={draw, minimum width=6mm, minimum height=6mm, inner sep=0pt, font=\scriptsize},
  box/.style={draw, rounded corners=2pt, minimum height=7mm, inner sep=2pt, font=\scriptsize, align=center},
  mlp/.style={draw, trapezium, trapezium angle=70, shape border rotate=270, minimum width=10mm, minimum height=9mm, inner sep=1pt, font=\scriptsize, fill=gray!25},
  par/.style={draw, rounded corners=4pt, minimum height=6mm, inner sep=2pt, font=\scriptsize},
  ar/.style={-{Stealth[length=1.6mm]}, thick},
  lab/.style={font=\scriptsize, inner sep=1pt}]
  % observation tape
  \foreach \i/\c in {1/bookaccent!30,2/bookexample!35,3/bookdef!35} \node[cell, fill=\c] (x\i) at (0.6*\i-0.6,0) {$x_{\i}$};
  \node[cell] (xd) at (1.8,0) {$\cdots$}; \node[cell, fill=bookpitfall!25] (xT) at (2.4,0) {$x_T$};
  % backward RNN
  \node[box, fill=bookexample!20, minimum width=20mm] (brnn) at (1.2,1.5) {backward LSTM\\(flip, LSTM, flip)};
  \draw[ar] (1.2,0.32) -- (brnn);
  \node[box, fill=bookkey!12] (bt) at (1.2,3.0) {$b_t$ (summary of $x_{t:T}$)};
  \draw[ar] (brnn) -- (bt);
  % combiner
  \node[box, fill=bookdef!25, minimum height=16mm, text width=12mm] (comb) at (4.3,3.5) {combiner};
  \node[box, fill=white] (zprev) at (3.0,5.0) {$\dot z_{t-1}$};
  \draw[ar] (bt.east) -- ([yshift=-4mm]comb.west);
  \draw[ar] (zprev.south) |- ([yshift=4mm]comb.west);
  % encoder
  \node[mlp] (enc) at (6.0,3.5) {MLP};
  \draw[ar] (comb) -- (enc);
  \node[par, fill=bookaccent!15] (mu) at (7.5,4.1) {$\mu_t$};
  \node[par, fill=bookaccent!15] (sg) at (7.5,2.9) {$\sigma_t$};
  \draw[ar] (enc.north east) -- (mu); \draw[ar] (enc.south east) -- (sg);
  \node[draw, circle, inner sep=0pt, minimum size=4.5mm, font=\scriptsize] (times) at (8.6,2.9) {$\odot$};
  \node[draw, circle, inner sep=0pt, minimum size=4.5mm, font=\scriptsize] (plus) at (9.4,3.5) {$+$};
  \node[par, fill=white] (eps) at (8.6,1.9) {$\varepsilon_t\sim\mathcal N(0, I)$};
  \draw[ar] (sg) -- (times); \draw[ar] (eps) -- (times);
  \draw[ar] (times) -- (plus); \draw[ar] (mu) -| (plus);
  \node[box, fill=white] (zt) at (10.4,3.5) {$\dot z_t$};
  \draw[ar] (plus) -- (zt);
  \draw[ar, densely dotted, gray!80!black] (zt.north) .. controls +(0,1.6) and +(1.5,0.6) .. node[lab, above, pos=0.2, fill=white]{next step} (zprev.east);
  % generative networks
  \node[mlp, fill=bookpitfall!15] (dec) at (10.4,1.6) {MLP};
  \draw[ar] (zt) -- (dec);
  \node[par, fill=bookpitfall!10] (mux) at (10.4,0.3) {$\mu^x_\theta(\dot z_t),\ \sigma^x_\theta(\dot z_t)$};
  \draw[ar] (dec) -- (mux);
  \node[mlp, fill=bookpitfall!15] (tr) at (5.6,5.9) {MLP};
  \draw[ar] (zprev.north) |- (tr.west);
  \node[par, fill=bookpitfall!10, anchor=west] (muz) at (6.4,5.9) {$\mu^z_\theta(\dot z_{t-1}),\ \sigma^z_\theta(\dot z_{t-1})$};
  \draw[ar] (tr) -- (muz);
  \node[lab, text=bookexample!70!black] at (1.2,4.6) {inference network};
  \node[lab, text=bookpitfall!80!black, anchor=west] at (6.4,6.55) {transition $p_\theta(z_t\mid z_{t-1})$};
  \node[lab, text=bookpitfall!80!black, anchor=north] at (10.4,-0.05) {decoder $p_\theta(x_t\mid z_t)$};
\end{tikzpicture}
\caption{One time step of the DKF in training. Inference network: the backward LSTM summarizes the
future, the combiner merges the summary with the previous sample, the encoder outputs the
parameters of $q_\phi(z_t\mid z_{t-1}, x_{t:T})$, and $\dot z_t$ is drawn by reparametrization. Generative networks (red):
the transition is evaluated at the sampled previous state, the decoder at the sampled current
state. Redrawn after the report of \texttt{CL-BZH/Deep-Kalman-Filters}.}
\label{dkf:fig:network}
\end{figure}

The forward pass of the repository's \texttt{DKF\_Base} makes the three stages explicit (abridged;
tensors have shape $(T, \text{batch}, \cdot)$):
\begin{lstlisting}[language=Python]
q_b_rnn_out = self.q_b_rnn(x)                  # b_1..b_T, one backward pass
z_t = z_0                                      # initial value, fixed or learned
self.combiner.reset()                          # use alpha_init at t = 1
for t in range(seq_len):
    combined = self.combiner(z_t, [q_b_rnn_out[t]])
    mu, logvar = self.encoder(combined)        # parameters of q(z_t | z_{t-1}, x_{t:T})
    z_t = self.sampler(mu, logvar)             # mu + exp(logvar/2) * eps
    self.z_mean[t], self.z_logvar[t], self.sampled_z[t] = mu, logvar, z_t
# transition p(z_t | z_{t-1}) at the sampled previous states, decoder at the samples
z_prev = torch.cat([z_0, self.sampled_z[:-1]], 0)
self.z_transition_mean, self.z_transition_logvar = self.p_z(z_prev)
self.x_hat, self.x_hat_logvar = self.decoder(self.sampled_z)
\end{lstlisting}
The loop over $t$ is unavoidable, because the combiner at step $t$ needs the sample of step $t-1$:
the trajectory $\dot z_{1:T}$ is drawn by ancestral sampling of $q_\phi$. The transition and the decoder, which only
read samples, are then applied to all time steps at once.

\section{The ELBO and its estimator}\label{dkf:sec:elbo}

The generative model \cref{dkf:eq:generative} and the inference model \cref{dkf:eq:q} are the two Markov
chains of \cref{em-algorithm:sec:mc-elbo}, with $y_t = x_t$ and with conditionals of $q$ that depend on the
observations. The derivation there gives
\[
  \mathrm{ELBO}(\theta, \phi) = \sum_{t=1}^T\E_{q_t}\bigl[\log p_\theta(x_t\mid z_t)\bigr]
  - \sum_{t=1}^T\E_{z_{t-1}\sim q_{t-1}}\Bigl[\mathrm{KL}\bigl(q_\phi(\cdot\mid z_{t-1}, x_{t:T})\,\Vert\,p_\theta(\cdot\mid z_{t-1})\bigr)\Bigr]
\]
(\cref{em-algorithm:eq:seq-elbo}), where $q_t$ is the marginal of $z_t$ under $q_\phi$; the outer expectation is trivial
for $t = 1$. The marginals have no closed form, but one ancestral trajectory $\dot z_{1:T}$, the one drawn
by the forward pass above, gives the unbiased estimator
\[
  \widehat{\mathrm{ELBO}} = \sum_{t=1}^T\log p_\theta(x_t\mid\dot z_t)
  - \sum_{t=1}^T\mathrm{KL}\bigl(q_\phi(\cdot\mid\dot z_{t-1}, x_{t:T})\,\Vert\,p_\theta(\cdot\mid\dot z_{t-1})\bigr)
\]
(\cref{em-algorithm:prop:unbiased}). Each KL term compares two Gaussians and has a closed form.

\begin{proposition}[KL between diagonal Gaussians]\label{dkf:prop:kl}
For $q = \mathcal N(\mu_q, \diag v_q)$ and $p = \mathcal N(\mu_p, \diag v_p)$ on $\R^d$,
\[
  \mathrm{KL}(q\,\Vert\,p) = \frac12\sum_{j=1}^d\Bigl(\log\frac{v_{p,j}}{v_{q,j}} + \frac{v_{q,j} + (\mu_{q,j} - \mu_{p,j})^2}{v_{p,j}} - 1\Bigr).
\]
\end{proposition}

\begin{proof}
Both densities are products over coordinates, so the KL is a sum over $j$, and it suffices to
treat $d = 1$. With $z\sim q$, $\E_q[(z - \mu_q)^2] = v_q$ and $\E_q[(z - \mu_p)^2] = v_q + (\mu_q - \mu_p)^2$, so
\[
  \E_q[\log q(z)] = -\tfrac12\log(2\pi v_q) - \tfrac12,\qquad
  \E_q[\log p(z)] = -\tfrac12\log(2\pi v_p) - \frac{v_q + (\mu_q - \mu_p)^2}{2v_p},
\]
and $\mathrm{KL}(q\Vert p) = \E_q[\log q(z)] - \E_q[\log p(z)]$ is the stated expression.
\end{proof}

\begin{example}[One KL term]\label{dkf:ex:kl}
In one dimension, $q = \mathcal N(0.5, 0.25)$ (standard deviation $0.5$) and $p = \mathcal N(0, 1)$:
$\mathrm{KL} = \frac12\bigl(\log4 + (0.25 + 0.25)/1 - 1\bigr) = \frac12(1.3863 - 0.5) = 0.4431$. The inference network pays this price for
moving the state half a unit away from what the dynamics predict and for being four times
more certain than they are.
\end{example}

The evaluation script of the repository implements the estimator with these two
functions, the reconstruction term with a Gaussian log-density:
\begin{lstlisting}[language=Python]
def gaussian_log_density(v, mean, logvar):     # summed over the feature dimension
    return -0.5 * (LOG_2PI + logvar + (v - mean) ** 2 / torch.exp(logvar)).sum(-1)

def gaussian_kl(mu_q, logvar_q, mu_p, logvar_p):
    return 0.5 * (logvar_p - logvar_q
                  + (torch.exp(logvar_q) + (mu_q - mu_p) ** 2) / torch.exp(logvar_p)
                  - 1.0).sum(-1)
\end{lstlisting}
One training step draws one trajectory per sequence of the mini-batch, evaluates
$-\widehat{\mathrm{ELBO}}/T$, backpropagates through the whole computation, the sampling steps included
thanks to the reparametrization, and updates $\theta$ and $\phi$ together. The backward LSTM is trained by
BPTT through the inference loop, like any other recurrent module.

\begin{remark}[The training loss of the repository]
The training loss (\texttt{Utils/loss\_fcts.py}) drops the factors $\tfrac12$ and the additive constants of
both terms: its reconstruction term is $\sum_{t,j}\bigl((x_{t,j} - \mu^x_j)^2/\sigma^{x\,2}_j + \log\sigma^{x\,2}_j\bigr)$ (or $\sum_t\norm{x_t - \mu^x}^2$ with the
\texttt{mse} option, a unit variance), and its KL term is $2\,\mathrm{KL}$ plus a constant, weighted by $\beta$. Both terms are
doubled, so for $\beta = 1$ the loss is $-2\,\widehat{\mathrm{ELBO}}/T$ up to an additive constant: the same minimizers, and
gradients that differ by the factor $2/T$ only. A value $\beta\neq1$ weights the KL term as in the $\beta$-VAE.
\end{remark}

\paragraph{Forecasting.} To predict the continuation of an observed sequence, infer a trajectory
$\dot z_{1:T}$ with the inference network, then run the transition forward from $\dot z_T$: draw $\dot z_{T+1}\sim p_\theta(\cdot\mid\dot z_T)$,
decode $\hat x_{T+1} = \mu^x_\theta(\dot z_{T+1})$, and repeat. The forecast depends on the observations only through $\dot z_T$.

\section{The autoregressive DKF: a transition that depends on the past}\label{dkf:sec:ar}

On sine waves, the vanilla DKF reconstructs the observed part of a sequence well but
forecasts it poorly: in the repository's experiments its prediction error is worse than that
of the constant prediction $0$. Its forecasts rely entirely on the learned transition
$p_\theta(z_t\mid z_{t-1})$, and nothing in training forces $z_t$ to carry the information needed to continue the
signal far ahead. The \emph{autoregressive DKF} (AR DKF), designed in the repository, gives the
dynamics direct access to the history: its transition is
\[
  p_\theta(z_t\mid z_{t-1}, x_{0:t-1}),
\]
a density for the next state that depends on the past observations, with a start value $x_0$
(zero or learned) so that the transition into $z_1$ is defined. This is exactly the situation of
\cref{sequence-models:sec:autoregressive}: a conditional distribution given a growing past, whose
past is summarized by a forward RNN run on the \emph{shifted} sequence $(x_0, x_1,\dots,x_{T-1})$. Its output
$h_{t-1}$ summarizes $x_{0:t-1}$, and an MLP maps $(z_{t-1}, h_{t-1})$ to the mean and log-variance of the transition:
\begin{lstlisting}[language=Python]
x_lagged = torch.cat((x_0, x[:-1]), 0)     # (x_0, x_1, ..., x_{T-1})
h_x_lagged = self.pz_f_rnn(x_lagged)       # position t: summary of x_{0:t-1}
z_lagged = torch.cat([z_0, self.sampled_z[:-1]], 0)
mean, logvar = self.p_z(z_lagged, h_x_lagged)   # p(z_t | z_{t-1}, x_{0:t-1})
\end{lstlisting}
The joint distribution becomes
\begin{equation}\label{dkf:eq:ar}
  p_\theta(x_{1:T}, z_{1:T}) = \prod_{t=1}^Tp_\theta(z_t\mid z_{t-1}, x_{0:t-1})\,p_\theta(x_t\mid z_t) .
\end{equation}
It is still a causal generative model: in the order $z_1, x_1, z_2, x_2,\dots$ each variable depends only on
variables generated before it, so ancestral sampling works and the product is normalized
(the argument of \cref{sequence-models:prop:normalized}). In the graph, the new factor adds an edge from
every past observation to $z_t$ (\cref{dkf:fig:ar}).

\begin{figure}[ht]
\centering
\begin{tikzpicture}[
  v/.style={draw, circle, minimum size=8.5mm, inner sep=0pt, font=\small},
  ob/.style={v, fill=gray!30},
  ar/.style={-{Stealth[length=1.8mm]}, thick},
  nw/.style={-{Stealth[length=1.8mm]}, thick, bookaccent}]
  \node[v, dashed] (z0) at (0,0) {$z_0$};
  \node[ob, dashed] (x0) at (0,-1.5) {$x_0$};
  \foreach \t in {1,2,3,4} {
    \node[v] (z\t) at (2.4*\t,0) {$z_{\t}$};
    \node[ob] (x\t) at (2.4*\t,-1.5) {$x_{\t}$};
    \draw[ar] (z\t) -- (x\t);
  }
  \draw[ar] (z0) -- (z1); \draw[ar] (z1) -- (z2); \draw[ar] (z2) -- (z3); \draw[ar] (z3) -- (z4);
  \draw[nw] (x0) -- (z1); \draw[nw] (x1) -- (z2); \draw[nw] (x2) -- (z3); \draw[nw] (x3) -- (z4);
  \draw[nw, densely dashed, bookaccent!60] (x1) .. controls +(1.6,-0.9) and +(-0.3,-1.6) .. (z3);
  \draw[nw, densely dashed, bookaccent!60] (x2) .. controls +(1.6,-0.9) and +(-0.3,-1.6) .. (z4);
\end{tikzpicture}
\caption{The AR DKF. The transition into $z_t$ depends on $x_{0:t-1}$ (blue): an edge goes from every past
observation to $z_t$; the edges from $x_{t-1}$ are drawn solid, two of the longer ones dashed, the others
omitted.}
\label{dkf:fig:ar}
\end{figure}

\paragraph{What the new edges change for inference.} Changing the generative model changes
its posterior, and therefore the inference network it needs.

\begin{proposition}[Posterior of the AR DKF]\label{dkf:prop:ar-posterior}
In the AR DKF,
\[
  p_\theta(z_{1:T}\mid x_{1:T}) = \prod_{t=1}^Tp_\theta(z_t\mid z_{t-1}, x_{1:T}),
\]
and the graph does not imply that $p_\theta(z_t\mid z_{t-1}, x_{1:T})$ is independent of the past observations $x_{1:t-1}$.
\end{proposition}

\begin{proof}
\emph{Earlier states drop out.} Let $s\leq t-2$ and take a path from $z_s$ to $z_t$, with all observations and
$z_{t-1}$ observed. If the path has an observation $x_r$ as an intermediate node, it is blocked there: $x_r$ has at
most one parent ($z_r$, none for $x_0$), so at most one of the two path edges at $x_r$ points into it, and $x_r$ is a chain or
fork node, which is observed. Otherwise the path uses latent nodes only, so it is the chain
$z_s\to\cdots\to z_{t-1}\to z_t$, blocked at the observed chain node $z_{t-1}$.

\emph{Past observations do not drop out.} For $t\geq2$, $x_{t-1}$ is a parent of $z_t$: the edge $x_{t-1}\to z_t$ is a path with no
intermediate node, which no conditioning set can block. By the second part of
\cref{dkf:thm:dsep}, the dependence holds for generic parameters.
\end{proof}

In the vanilla DKF the past observations influence $z_t$ only \emph{through} $z_{t-1}$, so knowing $z_{t-1}$ makes
them redundant. In the AR DKF they influence $z_t$ \emph{directly}, and they can never become
redundant: the inference network for $z_t$ must see the past as well as the future.

\paragraph{Three inference networks.} The repository implements three variants with the
generative model \cref{dkf:eq:ar}, which differ only in the inference network:
\begin{center}\small
\begin{tabular}{lp{0.5\textwidth}l}
\textbf{model} & \textbf{inference network} & \textbf{matches the posterior?}\\\hline
AR DKF & $q_\phi(z_t\mid z_{t-1}, x_{t:T})$: backward LSTM & no: ignores $x_{1:t-1}$\\
FAR DKF & $q_\phi(z_t\mid z_{t-1}, x_{1:T})$: backward LSTM and a second forward LSTM & yes\\
BAR DKF & $q_\phi(z_t\mid z_{t-1}, x_{1:T})$: one bidirectional LSTM & yes
\end{tabular}
\end{center}
In the FAR DKF (``forward autoregressive'') the combiner receives two summaries, $b_t$ of $x_{t:T}$ and
$\overrightarrow h_t$ of $x_{1:t}$, each with its own weight. In the BAR DKF (``bidirectional autoregressive'') the two
directions are merged inside a bidirectional RNN (\cref{sequence-models:sec:bidirectional}) before the
combiner. By \cref{dkf:prop:ar-posterior}, the AR DKF inference network cannot represent the
posterior of its own generative model, even with unlimited capacity: unless the true
posterior happens not to depend on the past observations, the gap
$\mathrm{KL}\bigl(q_\phi\,\Vert\,p_\theta(\cdot\mid x)\bigr)$ between the ELBO and the log-likelihood (\cref{em-algorithm:prop:gap}) cannot be
closed. The FAR and BAR variants remove this obstruction.

\paragraph{Prediction.} Forecasting works as for the DKF, except that each transition step
also needs the summary of the observations so far: after a forecast $\hat x_{T+s}$ is produced, it is
appended to the sequence and fed to the forward RNN, so the next transition is conditioned
on the model's own forecasts. The AR DKF is thus a genuinely autoregressive forecaster, with
the exposure-bias caveat of \cref{sequence-models:sec:autoregressive}. In the repository's
\texttt{predict}, the sequence fed to the forward RNN is a sliding window (each forecast is
appended and the oldest element dropped, and the RNN is re-run from the start of the
window), whereas in training the RNN always starts from $x_0$; the summaries seen at prediction
time are therefore computed slightly differently from those seen in training.

\paragraph{What the experiments show.} On the sine waves, the AR DKF forecasts far better
than the DKF: the prediction mean squared error is $0.876$ for the DKF and $0.0052$ for the AR DKF at
frequency $1$, and $0.571$ against $0.0094$ at frequency $2$ (always predicting $0$ gives $0.5$ for a unit
sine). On windows of the Lorenz attractor, over four training runs per model, the AR, FAR
and BAR DKF reconstruct equally well, none forecasts consistently better than the others,
and their estimated ELBO gaps are the same. The advantage that \cref{dkf:prop:ar-posterior}
suggests for the FAR and BAR DKF is not visible in these experiments: a structural argument
says what an inference network \emph{cannot} do, not how much the restriction costs in practice.

\begin{chaptersummary}
\begin{itemize}
  \item The DKF is a state-space model with neural transition and observation distributions;
        its posterior has no closed form and is approximated by an inference network.
  \item d-separation reads conditional independences from a Bayesian network: chains and forks
        block when their middle node is observed, colliders when neither they nor their
        descendants are.
  \item The DKF posterior factorizes as $\prod_tp(z_t\mid z_{t-1}, x_{t:T}, u_{t:T})$: the past is absorbed by $z_{t-1}$, the future
        observations matter. A backward LSTM summarizes $x_{t:T}$, a combiner merges the summary with
        $z_{t-1}$, and an encoder outputs a Gaussian sampled by reparametrization.
  \item The ELBO is estimated along the one ancestral trajectory of the forward pass, with
        closed-form Gaussian KL terms.
  \item The AR DKF makes the transition a density depending on the past, $p(z_t\mid z_{t-1}, x_{0:t-1})$, computed by a
        forward RNN on the shifted observations. Its posterior then depends on $x_{1:T}$, so the inference
        network needs a summary of the past too (FAR and BAR DKF).
\end{itemize}
\end{chaptersummary}

\section*{Exercises}
\addcontentsline{toc}{section}{Exercises}

\begin{exercise}\label{dkf:exr:1}
In the DKF without inputs, decide with d-separation whether (a) $x_1\indep x_3\mid z_2$; (b) $x_1\indep x_3$;
(c) $z_1\indep z_3\mid x_2$.
\solutionlink{dkf:exr:1}
\end{exercise}

\begin{exercise}\label{dkf:exr:2}
Starting from the product rule
$p(x_{1:T}, z_{1:T}) = \prod_tp(z_t\mid x_{1:t-1}, z_{1:t-1})\,p(x_t\mid x_{1:t-1}, z_{1:t})$, use d-separation to recover \cref{dkf:eq:generative}
(without inputs).
\solutionlink{dkf:exr:2}
\end{exercise}

\begin{exercise}\label{dkf:exr:3}
With inputs, could the inference factor for $z_t$ be $q_\phi(z_t\mid z_{t-1}, x_{t:T}, u_t)$, ignoring the future inputs?
Answer with the graph, and give an everyday example of the mechanism.
\solutionlink{dkf:exr:3}
\end{exercise}

\begin{exercise}\label{dkf:exr:4}
Compute $\mathrm{KL}(q\Vert p)$ and $\mathrm{KL}(p\Vert q)$ for $q = \mathcal N(0.5, 0.25)$ and $p = \mathcal N(0, 1)$, and show that
$\mathrm{KL}(q\Vert p) = 0$ if and only if $q = p$ for two diagonal Gaussians.
\solutionlink{dkf:exr:4}
\end{exercise}

\begin{exercise}\label{dkf:exr:5}
Is $z_t\indep x_{t-2}\mid z_{t-1}, x_{t-1}$ in the DKF? In the AR DKF?
\solutionlink{dkf:exr:5}
\end{exercise}

\begin{exercise}\label{dkf:exr:6}
List the recurrent networks of the AR, FAR and BAR DKF, with their direction, their input
sequence, what their state at time $t$ summarizes, and whether they belong to the generative
model or to the inference network.
\solutionlink{dkf:exr:6}
\end{exercise}
\end{notechapter}
